% ============================================================
% OP3_1_D1_blocks_v1.tex  (EXACT INSERTION BLOCKS for review;
% NOT yet implemented). Scope per GPT verdict on recon v2 + A0:
% OP3.1a/b/c decomposition near the OP3.1 item + one compact
% interface-requirements paragraph + status echo. Excluded:
% direct-method/coercive/minimiser scaffold; stochastic-
% quantisation route; finite-action claim; derives-mechanism
% language; any status upgrade.
% ============================================================

% ---------- ANCHOR 1 ----------
% In the OP3.1 decomposition item, REPLACE:
%   "(open; possible routes include configuration-space measure/saddle
%   selection and, more speculatively, stochastic-quantisation-type
%   relaxational formalisms);"
% WITH:
%   "(open; sharpened below into subproblems OP3.1a--c with recorded
%   interface requirements; mechanism open);"

% ---------- ANCHOR 2 ----------
% Insert AFTER the end of the decomposition list:
%   "\emph{OP3.4} --- relics of the ordering event (open observational
%    programme)."
% and BEFORE the paragraph "The second is: given that an ordered
% Lorentzian branch has formed, ...":

\paragraph{Sharpened decomposition of OP3.1 (Open; programme
sharpened).}
The ordering problem separates into three subproblems.
\emph{OP3.1a (bare-to-ordered gap).}
The bare micro-action is explicitly recorded
(eq.~\eqref{eq:ECT_action}, with field equation
\eqref{eq:ect_euclidean_fundamental}), but the gradient-ordered
sector is introduced by~P4 as a postulate: the current minimal
bare action does not by itself exhibit the gradient-ordered
sector, and the article records OP-grad as the missing bridge.
Exhibiting an ECT-native bridge --- an admissible $P(X)$-type
extension, a boundary-data mechanism of the recorded
boundary-value formulation, or another mechanism --- coincides
with completing OP-grad.
\emph{OP3.1b (EFT-data derivation).}
The task is to derive the ordered-branch EFT data $Z_u$,
$\mathcal C_n$, $W(u)$
of \eqref{eq:ordered_branch_euclidean_eft} by controlled expansion
around the ordered configuration --- hence $F$, $\omega$, $U$
within the recorded closure-function map
(\S\ref{sec:closure_function_mapping}) and its normalisation
freedoms; this subsumes the recorded matching problem~OP2.
Postulating $W(u)$ does not close OP3.1.
\emph{OP3.1c (selection as measure concentration).}
The task is to establish dominance of the ordered sector within
the recorded
$S_0$-weighted Euclidean extremal-action framework
(\S\ref{sec:qs_extremal_action}); ``selection'' here is a
statement about the Euclidean measure, not a process in time.

\paragraph{Interface requirements (recorded; no mechanism is
implied).}
Any candidate ordering mechanism must not presuppose emergent-time
or causal-cone structure prior to ordering; it must formulate
selection as Euclidean measure concentration rather than time
evolution; it must derive the ordered sector~(P4) and the EFT data
rather than assume them; it must hand $\rho_{\rm onset}$ to the
thermal-completion ledger only after an ordered Lorentzian branch
exists, with no energy statements at the pre-emergent Euclidean
level; and it must respect the recorded no-ghost expectations and
the no-phantom bound~\eqref{eq:w_phi}.
OP3.1 remains Open; interface requirements recorded / programme
sharpened.

% ---------- ANCHOR 3 ----------
% In the \statussummary Open list, REPLACE:
%   "(i)~first-principles ordering dynamics~(OP3.1); the branch"
% WITH:
%   "(i)~first-principles ordering dynamics~(OP3.1; interface
%   requirements recorded, subproblems OP3.1a--c separated above ---
%   programme sharpened, mechanism open); the branch"
